\clearpage\thispagestyle{empty}
\pdfbookmark[0]{Sources for Comet (15 February 1823)}{cometfeb1823-sources}
\begin{center}\textbf{Sources for Comet (15 February 1823)}\end{center}
\noindent This AI translation follows \textit{Comet} in \href{https://gdz.sub.uni-goettingen.de/id/PPN236007467}{\textit{Werke}, vol.~VI} (Göttingen, 1874), pp.~442--443. The first printing appeared under the heading \textit{Göttingen} in the \textit{Göttingische gelehrte Anzeigen}, no.~26, 15 February 1823, pp.~\href{https://images.p.sub.uni-goettingen.de/iiif/image/gdz:PPN385030444_1823_1:00000251/full/1600,/0/default.jpg}{249}, \href{https://images.p.sub.uni-goettingen.de/iiif/image/gdz:PPN385030444_1823_1:00000252/full/1600,/0/default.jpg}{250} and \href{https://images.p.sub.uni-goettingen.de/iiif/image/gdz:PPN385030444_1823_1:00000253/full/1600,/0/default.jpg}{251}. The scans are from the Göttingen State and University Library. The notice reports Rümker’s observations as communicated in a letter from Paramatta to Gauss. It refers to the journal’s \hyperlink{native-425}{report of 24 May 1819} on Encke’s discovery.
\par\smallskip\noindent The first printing spells Encke’s name \textit{Enke} and uses different time marks and decimal punctuation. It has no bracketed reference to p.~420 and no dotted leaders beside the two occultation times. The table and the occultation and longitude entries follow the collected edition, including the sparse units, the spelling \textit{Paramatta} and the sidereal hour 13 on 22 June, which is also printed in the journal. The report leaves open the meaning of \textit{mean} for the coordinates. Translation and typesetting produced by AI.
\clearpage
